\section{Circular Staking Prevention}
\label{sec:circular}

\subsection{Attack Description}

A circular staking attack attempts to amplify staking returns through leverage:
\begin{enumerate}
    \item Stake LUX $\to$ receive sLUX
    \item Use sLUX as collateral $\to$ borrow LUX
    \item Stake borrowed LUX $\to$ receive more sLUX
    \item Repeat
\end{enumerate}

If unchecked, this creates unbounded leverage, amplifying both returns and slashing losses.

\subsection{Prevention Mechanisms}

\begin{definition}[Anti-Circular Staking]
The LLS protocol prevents circular staking through:
\begin{enumerate}
    \item \textbf{sLUX is not accepted as staking collateral}: The StakingPool contract accepts only native LUX, not sLUX or wsLUX.
    \item \textbf{Deposit source verification}: Deposits from known lending protocol addresses are rate-limited.
    \item \textbf{Maximum leverage cap}: Even through indirect paths (borrow LUX from lending, stake, use sLUX as collateral), the maximum leverage is bounded by the lending LTV ratio.
\end{enumerate}
\end{definition}

\begin{theorem}[Leverage Bound]
With maximum LTV ratio $\ell \in (0, 1)$, the maximum leverage factor from recursive borrowing is:
\begin{equation}
    L_{\max} = \frac{1}{1 - \ell}
\end{equation}
For $\ell = 0.75$: $L_{\max} = 4$. A user starting with 1,000 LUX can stake at most 4,000 LUX equivalent through recursive lending.
\end{theorem}

\begin{proof}
Starting with principal $P$, after $k$ rounds of borrowing at LTV $\ell$:
\begin{equation}
    \text{Total staked} = P \sum_{i=0}^{k} \ell^i = P \cdot \frac{1 - \ell^{k+1}}{1 - \ell} \xrightarrow{k \to \infty} \frac{P}{1 - \ell}
\end{equation}
\end{proof}

\begin{proposition}[Bounded Slashing Amplification]
Under maximum leverage $L_{\max} = 4$, a 5\% slashing event causes at most $20\%$ loss on the user's original principal, which triggers liquidation (since $4 \times 0.05 = 0.20 > 0.075$ liquidation buffer), limiting systemic risk.
\end{proposition}

